\section*{References}\label{sec:references}
The sources are those named in the Lean files for the results that they prove or to which they
relate, as far as the paper cites them; the paper has not compared the statements it attributes to
them with the sources themselves. The Lean project's file \texttt{REFERENCES.md} says that ``only
some of the citations were compared with their sources'': an asterisk marks a work whose citation
it lists, for at least one Lean file, as ``compared, place confirmed'', and a dagger a work listed
only as ``place confirmed at second hand'' (``from a table of contents, a catalogue or a citation
in another work''); the other works are not reported there as compared.
\begingroup\small
\begin{list}{}{\setlength{\leftmargin}{1.5em}\setlength{\itemindent}{-1.5em}\setlength{\itemsep}{2pt}\setlength{\parsep}{0pt}}
\item Alon, N. (1999). Combinatorial Nullstellensatz.
\item Alp\"oge, L. (year not given). \url{https://alpo.ge/remark.pdf}. [Named, without title, in
  the \texttt{README} of the Lean project; cited as Alp\"oge's remark.]
\item Anderson (1991). [No title given.]
\item Ball, S., and Serra, O. (2009). Punctured combinatorial Nullstellens\"atze.
\item Barnet-Lamb, T., Gee, T., Geraghty, D., and Taylor, R. (2014). Potential automorphy and
  change of weight.
\item Behrend (1946). [No title given.] \emph{PNAS} \textbf{32}.
\item Bella\"iche and Chenevier (2009). Families of Galois representations and Selmer groups.
\item Benson and Doty (2009). Schur--Weyl duality over finite fields.
\item Bernstein (1984). $P$-invariant distributions on $\mathrm{GL}(N)$. \emph{LNM} \textbf{1041}.
\item Bertolini and Darmon (1996). [No title given.] \emph{Invent. Math.} \textbf{126}.
\item Bhatt, B., Morrow, M., and Scholze, P. (2018). Integral $p$-adic Hodge theory.
  \emph{Publ. IH\'ES} \textbf{128}.
\item Bloch, S., and Kato, K. (1990). $L$-functions and Tamagawa numbers of motives.
  \emph{Grothendieck Festschrift} I.
\item B\"ocherer, S. (1983). \"Uber die Fourier-Jacobi-Entwicklung Siegelscher Eisensteinreihen.
  \emph{Math. Z.} \textbf{183}.
\item Bombieri (1973)$^\dagger$. Counting points on curves over finite fields. \emph{S\'em.\ Bourbaki}
  430.
\item Bombieri, E., and Vaaler, J. (1983). On Siegel's lemma. \emph{Invent. Math.} \textbf{73}.
\item Borel, A. (2007). Automorphic forms on reductive groups (Park City lectures).
  \emph{IAS/Park City Math. Ser.} \textbf{12}, 7--39.
\item Bost (2001). Algebraic leaves of algebraic foliations over number fields.
\item Bourbaki, N. (1963)$^*$. Int\'egration VII.
\item Cais and Liu (2016). On $F$-crystalline representations. \emph{Doc. Math.} \textbf{21}.
\item Cassels, J. W. S. (1967). Global Fields. In Cassels--Fr\"ohlich, \emph{Algebraic Number
  Theory}, ch.~II.
\item Cassels and Fr\"ohlich (1967). Algebraic Number Theory.
\item Chenevier, G., and Harris, M. (2013). Construction of automorphic Galois representations, II.
  \emph{Camb. J. Math.} \textbf{1}.
\item Chida and Hsieh (2018). [No title given.] \emph{J. reine angew. Math.} \textbf{741}.
\item Clozel, L., Harris, M., and Taylor, R. (2008)$^*$. Automorphy for some $l$-adic lifts \dots\
  [title so abbreviated in the Lean project]. \emph{Publ. IHES} \textbf{108}.
\item Darmon, H., Diamond, F., and Taylor, R. (1995; revised 2007). Fermat's Last Theorem.
\item Dasgupta, Kakde, Silliman, and Wang (2023). The residually indistinguishable case of Ribet's
  method for $\mathrm{GL}_2$.
\item Dixon (1905). [No title given.]
\item Dotto and Thorne (2026). Residually reducible representations of rank 2.
\item Fakhruddin, Khare, and Patrikis (2021). Relative deformation theory, relative Selmer groups,
  and lifting irreducible Galois representations.
\item Forrester and Zhang (2021). Corank-1 projections and the randomised Horn problem.
\item Freitag, E. (1991). Singular Modular Forms and Theta Relations. \emph{Lecture Notes in Math.}
  \textbf{1487}, Springer.
\item Furusawa, M., and Morimoto, K. (2025). [No title given.] \emph{Math. Ann.} \textbf{391}.
\item Ghate and Jana (2025). Modular representations of $\mathrm{GL}_2(\mathbb{F}_q)$ using calculus.
\item Godement, R. (1957/58)$^\dagger$. S\'erie de Poincar\'e et Spitzenformen. \emph{S\'em.\ H.
  Cartan} \textbf{10}, exp.~10.
\item Gomez, R., and Zhu, C.-B. (2014). Local theta lifting of generalized Whittaker models.
  \emph{GAFA} \textbf{24}, 796--853.
\item Goodman and Wallach (2009)$^*$. Symmetry, Representations, and Invariants.
\item Green, J. A. (1955)$^*$. The characters of the finite general linear groups. \emph{Trans. AMS}
  \textbf{80}.
\item Green (1980). Polynomial representations of $\mathrm{GL}_n$.
\item Greither, C. (1992)$^*$. Class groups of abelian fields, and the main conjecture. \emph{Ann.
  Inst. Fourier} \textbf{42}.
\item Gross, B. (1999)$^\dagger$. Algebraic modular forms. \emph{Israel J. Math.} \textbf{113},
  61--93.
\item Harish-Chandra (1966). Discrete series for semisimple Lie groups II: explicit determination of
  the characters. \emph{Acta Math.} \textbf{116}. [Cited as Harish-Chandra, Discrete series II, 1966.]
\item Harris, Li, and Skinner (2005). The Rallis inner product formula and $p$-adic $L$-functions.
\item Heath-Brown, D. R. (1995). A mean value estimate for real character sums. \emph{Acta Arith.}
  \textbf{72}.
\item Howe, R. (1979). $\theta$-series and invariant theory. Corvallis 1979.
\item Imai and Koskivirta (2021). Automorphic vector bundles on the stack of $G$-zips.
\item Jutila, M. (1975). On mean values of Dirichlet polynomials with real characters. \emph{Acta
  Arith.} \textbf{27}.
\item Katz, N. (1978). $p$-adic $L$-functions for CM fields. \emph{Invent. Math.} \textbf{49}.
\item Katz (1981). Divisibilities, congruences, and Cartier duality.
\item Kisin, M. (2006)$^*$. Crystalline representations and $F$-crystals. \emph{Progr. Math.}
  \textbf{253}, 459--496.
\item Kneser, M. (1969). Lectures on Galois cohomology of classical groups. \emph{TIFR}
  \textbf{47}.
\item Kriz and Nordentoft (2023). Horizontal $p$-adic $L$-functions.
\item Kudla, S. (1986). On the local theta-correspondence. \emph{Invent. Math.} \textbf{83}.
\item Labesse, J.-P. (2011)$^*$. Changement de base CM et s\'eries discr\`etes. In \emph{On the
  stabilization of the trace formula}.
\item Lagarias, J. C., Lenstra, H. W., and Schnorr, C. P. (1990). Korkin--Zolotarev bases and
  successive minima of a lattice and its reciprocal lattice. \emph{Combinatorica} \textbf{10}.
\item Landherr, W. (1936). \"Aquivalenz Hermitescher Formen \"uber einem beliebigen algebraischen
  Zahlk\"orper. \emph{Abh. Math. Sem. Hamburg} \textbf{11}, 245--248.
\item Lang, S. (year not given). Algebraic Number Theory.
\item Laurent (1994). Linear forms in two logarithms and interpolation determinants.
\item Li, J.-S. (1992). Non-vanishing theorems for the cohomology of certain arithmetic quotients.
  \emph{Crelle} \textbf{428}.
\item Lubin and Tate (year not given). Formal complex multiplication in local fields.
\item Mazur, B., and Wiles, A. (1984). Class fields of abelian extensions of $\mathbb{Q}$.
  \emph{Invent. Math.} \textbf{76}.
\item Mazur and Rubin (2018). Diophantine stability.
\item Milne, J. S. (2006). Arithmetic Duality Theorems, 2nd ed.
\item Milne, J. S. (2020)$^*$. Class Field Theory, v4.03.
\item Milne, J. (2020)$^\dagger$. Complex Multiplication.
\item Montgomery, H. L., and Vaughan, R. C. (2007)$^*$. Multiplicative number theory I. Classical
  theory.
\item Mumford (year not given). Tata lectures on theta I.
\item Poor, C., and Yuen, D. (2000). Linear dependence among Siegel modular forms. \emph{Math. Ann.}
  \textbf{318}, 205--234.
\item Rallis, S. (1982). Langlands' functoriality and the Weil representation. \emph{Amer. J. Math.}
  \textbf{104}.
\item Rallis, S. (1984). Injectivity properties of liftings associated to Weil representations.
  \emph{Compositio} \textbf{52}. [Cited with its title, to tell it from the other work of 1984.]
\item Rallis, S. (1984). On the Howe duality conjecture. \emph{Compositio Math.} \textbf{51},
  333--399. [Cited with its title, to tell it from the other work of 1984.]
\item Rallis, S. (1987). $L$-functions and the oscillator representation. \emph{LNM} \textbf{1245}.
\item Ribet (1976). [No title given.]
\item Rouquier, R. (1996). Caract\'erisation des caract\`eres et pseudo-caract\`eres. \emph{J.
  Algebra} \textbf{180}.
\item Satake, I. (1957/58)$^*$. Caract\'erisation de l'espace des Spitzenformen. \emph{S\'em.\ H.
  Cartan} \textbf{10}, exp.~9 bis.
\item Selberg, A., and Chowla, S. (1967)$^\dagger$. On Epstein's zeta-function. \emph{J. reine angew.
  Math.} \textbf{227}, 86--110.
\item Serre, J.-P. (1997). Galois Cohomology.
\item Shapira (2006). Behrend-type constructions for sets of linear equations. \emph{Acta Arith.}
\item Shimura, G. (1971)$^\dagger$. Introduction to the Arithmetic Theory of Automorphic Functions.
\item Shimura, G. (1975). On some arithmetic properties of modular forms of one and several
  variables. \emph{Ann. of Math.} \textbf{102}.
\item Shin, S. W. (year not given). Galois representations arising from some compact Shimura
  varieties.
\item Siegel, C. L. (1935). \"Uber die analytische Theorie der quadratischen Formen. \emph{Ann. of
  Math.} \textbf{36}.
\item Silverman, J. (year not given)$^*$. The Arithmetic of Elliptic Curves.
\item Tadi\'c (1986). Classification of unitary representations in irreducible representations of
  general linear group (non-archimedean case). \emph{Ann. ENS} \textbf{19}.
\item Tate, J. (1950)$^*$. Fourier analysis in number fields and Hecke's zeta-functions. Thesis; in
  Cassels--Fr\"ohlich, \emph{Algebraic Number Theory}, ch.~XV.
\item Taylor, R. (1991). Galois representations associated to Siegel modular forms of low weight.
  \emph{Duke Math. J.} \textbf{63}.
\item Taylor, R., and Wiles, A. (1995). Ring-theoretic properties of certain Hecke algebras.
  \emph{Ann. Math.} \textbf{141}.
\item Urban (2001). Selmer groups and the Eisenstein--Klingen ideal.
\item Walls, P. (2015). The theta correspondence, periods of automorphic forms and special values of
  standard automorphic $L$-functions. arXiv 1501.01565.
\item Wang Peng Jun, B. (2025)$^*$. Global theta lifting and automorphic periods associated to
  nilpotent orbits. \emph{J. Number Theory} \textbf{271}, 122--149.
\item Weil, A. (1965)$^*$. Sur la formule de Siegel dans la th\'eorie des groupes classiques.
  \emph{Acta Math.} \textbf{113}.
\item Wiles (1990). The Iwasawa conjecture for totally real fields.
\item Wu, C. (2013). Irreducibility of theta lifting for unitary groups. \emph{J. Number Theory}
  \textbf{133}.
\item Yang, T. (1997). Theta liftings and Hecke $L$-functions. \emph{J. reine angew. Math.}
  \textbf{485}.
\end{list}\endgroup
